\section{The problem and the approach}\label{ms:introduction}

The central limit theorem describes the shape of a long sum. The
Berry--Esseen inequality asks how far that description can be from the
truth at a specified sample size. Even when the summands have the same
distribution, the best universal constant is a delicate question: the
largest errors come from the interaction of skewness with the jumps of a
lattice distribution. A proof of the sharp inequality must control both
effects uniformly in the law and in the sample size.

Let $\mathcal P$ denote the probability laws on $\R$ satisfying
$\E X=0$, $\E X^2=1$, and $\rho(P):=\E|X|^3<\infty$. For independent
$X_1,\ldots,X_n$ with common law $P$, write $S_n=X_1+\cdots+X_n$ and
\begin{align*}
 \Delta_n(P)&=\sup_{H\in\R}
       |\Pp_P(S_n\le H)-\Phi(H/\sqrt n)|,\\
 R_n(P)&=\frac{\sqrt n\,\Delta_n(P)}{\rho(P)},\qquad
 C_n=\sup_{P\in\mathcal P}R_n(P),\qquad C^*=\sup_{n\ge1}C_n.
\end{align*}
Here $\Phi$ and $\phi$ are the standard normal distribution function and
density. The conjecture is
\begin{equation}\label{ms:conjecture}
 C^*=\CE,\qquad
 \CE=\frac{3+\sqrt{10}}{6\sqrt{2\pi}}=0.409732\ldots.
\end{equation}
Throughout, $H$ is a threshold for the \emph{raw sum} $S_n$ and
$H/\sqrt n$ its standardized version. All terminating decimals
in numerical inequalities mean exact rational numbers.

The normalization costs no generality. If $Z$ has mean $\mu$, variance
$\varsigma^2>0$, and finite third absolute central moment $\beta_3$, apply the
definitions to $X=(Z-\mu)/\varsigma$. Then $\rho=\beta_3/\varsigma^3$ and
\[
 \Pp\left(\frac{\sum_{i=1}^n Z_i-n\mu}{\varsigma\sqrt n}\le x\right)
 =\Pp(S_n\le\sqrt n\,x).
\]
Thus $C_n\le\CE$ is precisely the usual sharp inequality at index $n$.

\subsection{From Berry and Esseen to the sharp-constant problem}\label{ms:history}

Berry~\cite{Berry1941} and Esseen~\cite{Esseen1942,Esseen1945}
established quantitative normal approximation of order $n^{-1/2}$ under
a finite third absolute moment. Their results initiated the search for
the smallest constant in a bound that is uniform over the summand law.
Esseen's later moment analysis~\cite{Esseen1956} identified the lower
constant in~\eqref{ms:conjecture}. The conjecture that this lower constant
is also the universal upper constant is associated with Zolotarev's work
of 1966--1967~\cite{Zolotarev1966,Zolotarev1967}.

The two-point example explains the value of the constant. Put
\[
 p_*=(4-\sqrt{10})/2,\qquad q_*=1-p_*,\qquad
 X=(B-p_*)/\sqrt{p_*q_*},\quad B\sim\operatorname{Bernoulli}(p_*).
\]
For this fixed law, $\rho=(p_*^2+q_*^2)/\sqrt{p_*q_*}$ and the
classical lattice asymptotic gives
\begin{equation}\label{ms:bernoulli-limit}
 \lim_{n\to\infty}R_n(P)
 =\frac{3+q_*-p_*}{6\sqrt{2\pi}(p_*^2+q_*^2)}=\CE.
\end{equation}
Indeed $p_*^2+q_*^2=10-3\sqrt{10}$ and
$(3+q_*-p_*)/(p_*^2+q_*^2)=3+\sqrt{10}$.
Consequently every constant smaller than $\CE$ fails for a sufficiently
large sum of this one fixed Bernoulli law. This proves $C^*\ge\CE$;
it supplies no upper bound for other laws. Hipp and Mattner settled the symmetric binomial case
\cite{HippMattner2008}; Schulz's dissertation~\cite{Schulz2016} proves
the sharp upper bound for the entire binomial family at every sample size. We use Schulz's theorem for the noiseless case in
Section~\ref{ms:counts}.

Upper estimates developed through increasingly precise smoothing and
characteristic-function inequalities. Zolotarev's work and van Beek's
Fourier analysis~\cite{Zolotarev1967,vanBeek1972} are early milestones;
Prawitz's one-sided smoothing kernel~\cite{Prawitz1972} remains useful in
modern computations. Shiganov obtained the iid bound $0.7655$ in work originally published in
1982~\cite{Shiganov1986,Tyurin2009}. Subsequent improvements include Shevtsova's bound
$0.7056$~\cite{Shevtsova2007}, the bound $0.5129$ of Korolev and
Shevtsova~\cite{KorolevShevtsova2010}, Tyurin's bound $0.4785$ obtained
using zero-bias estimates and characteristic-function bounds~\cite{Tyurin2010}, and Shevtsova's bounds
$0.4748$ and $0.4690$~\cite{Shevtsova2011,Shevtsova2013}.
Their structural forms are more informative than the constants alone.
For example, the inequalities
\begin{equation}\label{ms:structural-classical}
 \sqrt n\,\Delta_n(P)\le
 \min\{0.4690\rho,\ 0.3322(\rho+0.429),\
                          0.3031(\rho+0.646)\}
\end{equation}
hold for every $P\in\mathcal P$ and $n\ge1$~\cite{Shevtsova2013}.
If $R_n(P)$ is near $\CE$, they bound $\rho$ from above. This supplies the
first quantitative moment cap for the support calculation.

\subsection{Recent manuscripts, repositories, and formalization}

Several recent projects combine analytic inequalities with extensive
computation. Xiao and Li report a global iid upper bound of $0.4395$;
Lin reports $0.423$ using a refinement of their analytic framework
\cite{xiao2026berryesseen,lin2026berryesseen}. Both supply Lean
developments and exact finite certificates. Tao and contributors' \emph{Optimization Constants in Mathematics},
especially Problem~19a, records this work and its verification status
\cite{tao2026berryesseenregistry}.

He and Cheng's preprint~\cite{hecheng2026eventually} reports the sharp
inequality for every
\[
 n\ge 2\lceil\exp(10^{17})\rceil,
\]
with a threshold independent of the summand law. This uniform quantifier
is essential. A theorem for each fixed law, with a law-dependent
threshold, would not exclude a sequence of counterexamples whose laws
change with $n$. A finite but astronomically large range of $n$ remains.

Bell's collection treats independent, not necessarily identically
distributed, summands. It reports a universal upper bound below
$0.44988$, results for finite arrays, and sharp bounds when the Lyapunov
ratio is small \cite{BellIndependent2026,BellFinite2026,BellSmall2026};
its structural papers study extremizers in that
setting~\cite{BellExtremizers2026,BellStructural2026}.

In Lean, Yan's \emph{ProbabilityApproximation} formalizes classical
nonuniform and multivariate normal-approximation
theorems~\cite{Yan2026ProbabilityApproximation}, and the
\emph{BerryEsseenLean} project formalizes an eventual-sharpness framework accompanying an earlier version of He
and Cheng, relative to explicitly listed external mathematical inputs~\cite{BerryEsseenLean2026}.

\subsection{The present approach and its scope}\label{ms:scope}

The analytic starting point is a maximizing law at a fixed sample size.
Its first variation gives a contact inequality at every support point.
We integrate that inequality against the density of the zero-bias
transform, introduced by Goldstein and Reinert~\cite{GoldsteinReinert1997}.
The resulting identity connects the normal-approximation error to a
nonnegative moment defect and a Fourier remainder. The defect singles out
Esseen's two-point law; its quantitative form controls conditional noise
and the displacement of the two parts of the support at the same time.
The sign decomposition then turns the original sum into an exact binomial
mixture. Related zero-bias methods appear in Tyurin's work and, for
independent summands, in Bell's transport paper~\cite{Tyurin2010,BellTransport2026}.

The reductions developed here address four parts of the argument.
The one-step recurrence of Theorem~\ref{ms:one-step} controls the first
index $n\ge42$ with $C_n\ge\CE$, assuming a sharp base through $41$;
Lemma~\ref{ms:first-failure-entry} then gives zero predecessor loss and
$0\le n(C_n-\CE)<3003/10250$. Proposition~\ref{ms:04:three-bounds}
uses this zero loss to prove support bounds by boundary estimates and
monotonicity. Retaining $n(C_n-\CE)\rho$ in the conserved identity
sharpens the joint bound on noise and endpoint displacement
(Corollary~\ref{ms:conserved:phase-resource} and
Lemma~\ref{ms:06:charge}). Finally, Theorem~\ref{ms:curvature-thm}
replaces the curvature maximization by at most three explicit candidates.
The support and resource statements keep their stated moment and
Fourier hypotheses.

The all-index theorem is now established unconditionally. Theorem~\ref{ms:conditional-global}
shows first that a sharp base $C_j\le\CE$ for $1\le j\le41$,
together with the equality-inclusive local exclusion of
Assumption~\ref{ms:local42}, implies $C_n<\CE$ for $n\ge42$ and $C^*=\CE$;
and second, that combining the analytic reductions proved here and in the
appendices with the complete thirteen-role finite and effective-tail
certificate family verified end-to-end (\texttt{full\_pass: "ESTABLISHED"})
proves
\[
 C_n\le\CE\quad(n\ge1),\qquad C^*=\CE=\frac{3+\sqrt{10}}{6\sqrt{2\pi}}.
\]
Establishing the all-index bound combines these analytic implications with
complete finite certificates verified across all thirteen index roles
($1\le n\le48$ and $n\ge49$). Section~\ref{ms:certificates}
records the global theorem, the thirteen-role partition, and the verified
\texttt{full\_pass} execution status in one place.

An earlier, AI-assisted working draft is available on
Hexagon~\cite{rodziewiczkolodziejek2026momentdefects}; the present paper
reorganizes that material and develops the reductions above. None of the
bounds reported in the recent works cited above is used here.
